Con el propósito de cumplir el primer objetivo específico —identificar las dificultades conceptuales que enfrentan los estudiantes de la asignatura \textit{Fundamentos de Interpretación y Compilación de Lenguajes de Programación}— se diseñó y aplicó un proceso de recolección de información que combina dos instrumentos complementarios: una encuesta dirigida a estudiantes que han cursado la asignatura y una encuesta dirigida al cuerpo docente responsable de impartirla. Los resultados de ambos instrumentos se analizan de forma integrada para identificar los conceptos con mayor nivel de dificultad, los factores que los explican y las implicaciones para el diseño del prototipo web propuesto.

\section{Encuesta a estudiantes}

\subsection{Diseño del instrumento}

La encuesta para estudiantes fue diseñada con el objetivo de caracterizar las dificultades conceptuales percibidas durante el cursado de la asignatura, así como la percepción sobre las herramientas y recursos disponibles. El cuestionario fue estructurado en cuatro bloques temáticos:

\begin{enumerate}
    \item \textbf{Bloque I - Perfil académico:} semestre en que cursó la asignatura, nota final obtenida y si la ha repetido.
    \item \textbf{Bloque II - Dificultades conceptuales:} nivel de dificultad percibido en cada uno de los temas centrales de la asignatura, mediante escala Likert de 5 niveles (1: sin dificultad, 5: dificultad muy alta).
    \item \textbf{Bloque III - Herramientas y recursos:} percepción sobre la utilidad de Racket y la librería EOPL como herramientas de apoyo; dificultad para encontrar recursos externos complementarios.
    \item \textbf{Bloque IV - Recursos adicionales:} valoración sobre si recursos interactivos o visualizaciones habrían mejorado la comprensión de los conceptos.
\end{enumerate}

El instrumento completo se puede consultar en el Anexo~\ref{annex:encuesta-estudiantes}.

\subsection{Población y muestra}

La encuesta fue respondida por 27 estudiantes que cursaron la asignatura en diferentes periodos académicos del programa de Ingeniería de Sistemas de la Universidad del Valle, sede Tuluá. Dado que el objetivo no era realizar inferencia estadística sino caracterizar cualitativamente las dificultades percibidas, la muestra se considera suficiente para identificar tendencias representativas del contexto local.

\subsection{Resultados}

A continuación se presentan los resultados agrupados por bloque temático.

\subsubsection{Bloque II — Dificultades conceptuales por tema}

\begin{longtable}{|p{7cm}|p{2cm}|p{5.5cm}|}
    \hline
    \textbf{Concepto} & \textbf{Promedio} & \textbf{Observación principal} \\
    \hline
    \endhead
    Definición de gramáticas libres de contexto (BNF/EBNF) & 3.7 / 5 & Dificultad moderada-alta; se confunden reglas de producción con expresiones regulares. \\
    \hline
    Análisis sintáctico (\textit{parsing}) y algoritmos LL/LR & 4.1 / 5 & Dificultad alta; los estudiantes no logran asociar la gramática con el árbol resultante. \\
    \hline
    Árboles de sintaxis abstracta (AST) & 4.3 / 5 & Dificultad muy alta; concepto abstracto sin representación visual de apoyo. \\
    \hline
    Ambientes de ejecución y alcance léxico & 4.2 / 5 & Dificultad muy alta; no se visualiza cómo cambia el ambiente durante la evaluación. \\
    \hline
    Clausuras y funciones de orden superior & 4.0 / 5 & Dificultad alta; los estudiantes comprenden la teoría pero no su comportamiento en ejecución. \\
    \hline
    Recursión en contextos de interpretación & 3.5 / 5 & Dificultad moderada; mayor familiaridad por su presencia en otros cursos. \\
    \hline
    Definiciones y esquemas dirigidos por sintaxis & 3.9 / 5 & Dificultad alta; los estudiantes no conectan la gramática con la semántica del programa. \\
    \hline
    \caption{Dificultad percibida por los estudiantes en los conceptos centrales de la asignatura FLP.\\
    \scriptsize \textbf{Fuente:} Elaboración propia. Escala Likert 1–5 (n = 27).}
    \label{tab:dificultad-conceptos}
\end{longtable}

Los tres conceptos con mayor dificultad promedio son: (1)~árboles de sintaxis abstracta (4.3), (2)~ambientes de ejecución y alcance léxico (4.2) y (3)~análisis sintáctico (4.1). Estos resultados coinciden con los reportados en la literatura especializada, donde se señala que la visualización de estructuras internas constituye una de las principales barreras pedagógicas en la enseñanza de compiladores e intérpretes \cite{Stamenković2024AWE, delvado2023itt}.

\subsubsection{Bloque III — Percepción sobre herramientas actuales}

\begin{longtable}{|p{8.5cm}|p{2.2cm}|p{3.8cm}|}
    \hline
    \textbf{Ítem} & \textbf{Respuesta afirmativa} & \textbf{Observación} \\
    \hline
    \endhead
    Racket fue una herramienta útil para comprender los conceptos de la asignatura & 29.6\% & La mayoría no lo percibió como apoyo didáctico. \\
    \hline
    La librería EOPL facilitó la comprensión práctica del material & 22.2\% & Muy baja percepción de utilidad pedagógica. \\
    \hline
    Encontré recursos externos complementarios con facilidad & 25.9\% & Mayoría tuvo dificultad para hallar material de apoyo. \\
    \hline
    Sentí seguridad al aplicar los conceptos teóricos en los proyectos & 33.3\% & Dos terceras partes expresaron incertidumbre o dificultad de aplicación. \\
    \hline
    \caption{Percepción de los estudiantes sobre las herramientas y recursos actuales de la asignatura.\\
    \scriptsize \textbf{Fuente:} Elaboración propia (n = 27).}
    \label{tab:percepcion-herramientas}
\end{longtable}

Los datos evidencian que Racket y la librería EOPL, aunque son herramientas técnicamente adecuadas, no son percibidas como recursos que faciliten la comprensión pedagógica de los conceptos. Este hallazgo está alineado con lo señalado por Spigariol et al.\ \cite{Spigariol2023Aprendiendo}, quienes indican que los intérpretes reales no exponen con claridad los procesos internos que llevan a la ejecución del código, limitando su aprovechamiento con fines didácticos.

\subsubsection{Bloque IV — Valoración de recursos adicionales}

El resultado más significativo de la encuesta fue que \textbf{la totalidad de los 27 participantes} (100\%) indicó que contar con recursos adicionales —en particular, herramientas visuales e interactivas que permitieran explorar la construcción del AST y los cambios en el ambiente de ejecución— habría mejorado de forma notable su proceso de aprendizaje. Este consenso constituye una motivación directa para el desarrollo del prototipo propuesto.

\section{Encuesta a docentes}

\subsection{Diseño del instrumento}

Con el fin de complementar la perspectiva de los estudiantes con la visión del cuerpo docente, se diseñó una encuesta semiestructurada dirigida a los profesores que han dictado la asignatura. El instrumento contempla cuatro dimensiones:

\begin{enumerate}
    \item \textbf{Perfil docente:} experiencia impartiendo la asignatura, enfoque metodológico habitual.
    \item \textbf{Dificultades observadas:} conceptos en los que los docentes identifican mayores falencias en sus estudiantes.
    \item \textbf{Factores asociados:} causas que, a juicio del docente, explican dichas dificultades.
    \item \textbf{Herramientas didácticas:} percepción sobre los recursos actuales y apertura hacia herramientas interactivas.
\end{enumerate}

El instrumento completo se puede consultar en el Anexo~\ref{annex:encuesta-docentes}.

\subsection{Población y muestra}

La encuesta fue respondida por 3 docentes con experiencia en la impartición de la asignatura \textit{Fundamentos de Interpretación y Compilación de Lenguajes de Programación} en la Universidad del Valle. Dado el tamaño reducido del cuerpo docente que dicta la asignatura en la sede Tuluá, se consideró a todos los docentes disponibles.

\subsection{Resultados}

A continuación se presentan los resultados agrupados por dimensión.

\subsubsection{Dificultades observadas en estudiantes}

\begin{longtable}{|p{7cm}|p{2.5cm}|p{5cm}|}
    \hline
    \textbf{Concepto / dificultad observada} & \textbf{Docentes que lo señalan} & \textbf{Observación} \\
    \hline
    \endhead
    Comprensión de gramáticas formales y notación BNF & 3/3 & Totalidad lo menciona como punto crítico. \\
    \hline
    Construcción e interpretación de árboles de sintaxis abstracta (AST) & 3/3 & Concepto que genera mayor confusión según los docentes. \\
    \hline
    Gestión y visualización de ambientes de ejecución & 3/3 & Los estudiantes no logran trazar mentalmente los cambios en el ambiente. \\
    \hline
    Alcance léxico y clausuras & 2/3 & Dificultad alta, especialmente al combinar con recursión. \\
    \hline
    Definiciones dirigidas por sintaxis y semántica operacional & 2/3 & Abstracción difícil de asimilar sin representaciones visuales. \\
    \hline
    \caption{Dificultades observadas por los docentes en los estudiantes de la asignatura FLP.\\
    \scriptsize \textbf{Fuente:} Elaboración propia (n = 3 docentes).}
    \label{tab:dificultad-docentes}
\end{longtable}

\subsubsection{Factores asociados a las dificultades}

Los tres docentes coincidieron en que una causa determinante de las dificultades es que los estudiantes llegan a la asignatura \textbf{con falencias en cursos previos}, particularmente en matemáticas discretas, teoría de autómatas y estructuras de datos. Esta situación limita la capacidad de los estudiantes para razonar sobre estructuras formales, lo que incrementa la brecha entre los fundamentos necesarios y los contenidos de la asignatura.

Adicionalmente, los docentes señalaron:

\begin{itemize}
    \item La alta carga cognitiva inherente a los contenidos dificulta el aprendizaje con los materiales tradicionales.
    \item Racket y EOPL ofrecen una base sólida, pero su curva de aprendizaje técnico consume tiempo que podría dedicarse a los conceptos de fondo.
    \item La falta de recursos visuales e interactivos reduce la capacidad de los estudiantes para construir modelos mentales de los procesos de interpretación.
\end{itemize}

\subsubsection{Percepción sobre herramientas interactivas}

Los tres docentes consideraron que una herramienta web interactiva que permitiera definir gramáticas sencillas, visualizar los árboles de sintaxis abstracta resultantes y observar cómo evoluciona el ambiente de ejecución durante la evaluación de un programa, podría reducir significativamente la brecha entre la comprensión teórica y la aplicación práctica. Señalaron como características deseables: acceso sin instalación, ejemplos predefinidos incrementales y retroalimentación inmediata sobre los cambios generados al modificar una expresión.

\section{Síntesis: Conceptos prioritarios para el prototipo}

\label{sec:conceptos-prioritarios}

A partir del análisis integrado de las encuestas a estudiantes y docentes, se identificaron los \textbf{cinco conceptos con mayor nivel de dificultad y mayor relevancia pedagógica} para orientar el diseño del prototipo web:

\begin{enumerate}
    \item \textbf{Gramáticas libres de contexto:} los estudiantes confunden reglas de producción con expresiones regulares y tienen dificultad para derivar cadenas a mano. El prototipo debe permitir definir gramáticas sencillas e interactuar con ellas en tiempo real.

    \item \textbf{Análisis sintáctico (parsing):} la relación entre una gramática y el proceso que produce un árbol a partir de una cadena es poco comprensible sin apoyo visual. El prototipo debe mostrar paso a paso cómo se construye el árbol a partir de la gramática definida.

    \item \textbf{Árboles de sintaxis abstracta (AST):} al ser una estructura puramente conceptual durante la enseñanza, la falta de representación visual es el factor que más contribuye a su incomprensión. El prototipo debe renderizar gráficamente el AST y permitir explorar sus nodos.

    \item \textbf{Ambientes de ejecución y alcance léxico:} los estudiantes no logran trazar mentalmente cómo el ambiente se extiende o se encadena durante la evaluación de expresiones. El prototipo debe representar visualmente los cambios en el ambiente a medida que se evalúa el programa.

    \item \textbf{Clausuras y funciones de orden superior:} aunque los estudiantes conocen el concepto, no comprenden su comportamiento al momento de la aplicación. La representación del ambiente capturado por una clausura, accesible durante la ejecución, puede materializar este concepto abstracto.
\end{enumerate}

Estos cinco conceptos orientarán la selección de ejemplos pedagógicos, los casos de uso y los mecanismos de visualización del prototipo, y servirán como base para la biblioteca de casos predefinidos que se incluirá en la herramienta.

\section{Implicaciones para el diseño del prototipo}

Los resultados del proceso de identificación de dificultades permiten derivar las siguientes implicaciones directas para el diseño del prototipo web:

\begin{itemize}
    \item \textbf{Visualización de AST como característica central:} dado que es el concepto con mayor dificultad percibida y el que más señalan docentes y estudiantes, la visualización gráfica e interactiva del árbol de sintaxis abstracta debe constituir el elemento principal de la interfaz.

    \item \textbf{Representación explícita del ambiente:} el ambiente de ejecución debe mostrarse de forma dinámica, reflejando cómo cambia con cada evaluación de expresión, para que el estudiante pueda rastrear el flujo de valores e identificadores.

    \item \textbf{Gramáticas configurables y ejemplos predefinidos:} dado que los estudiantes llegan con falencias previas, el prototipo debe incluir gramáticas de ejemplo incrementales (de menor a mayor complejidad) que permitan al estudiante orientarse antes de definir las propias.

    \item \textbf{Retroalimentación inmediata:} la herramienta debe proporcionar respuesta visual inmediata ante cada modificación del programa o la gramática, reduciendo así la carga cognitiva extrínseca \cite{SWELLER1988257, sweller2024cognitive}.

    \item \textbf{Accesibilidad sin instalación:} dado que Racket y EOPL presentan barreras técnicas de configuración, el prototipo debe funcionar en cualquier navegador web moderno sin requisitos de instalación, para que los estudiantes puedan acceder desde el primer momento.
\end{itemize}
